% ============================================================
% discussion.tex
% Sección: Discusión
% Tesis: Impacto del Gasto en Limpieza Pública sobre
%        Gestión de Residuos Sólidos Municipales, Perú 2015-2019
% Autor: Dante Barreto Gamarra (UNAS)
% ============================================================

\section{Discusión}
\label{sec:discusion}

\subsection{Magnitud de los Efectos y Benchmarks Comparativos}
\label{sec:benchmarks}

La elasticidad estimada de $+0{,}879$ entre el gasto municipal y la cantidad diaria de residuos recolectados se sitúa en el rango alto de los efectos documentados en la literatura sobre costos y eficiencia de la gestión de residuos sólidos. \citet{Bel2010}, en su análisis de municipios catalanes, documentan que las economías de escala en la gestión de residuos se agotan a partir de poblaciones superiores a $20{,}000$ habitantes, lo que implica elasticidades gasto-producción decrecientes en municipios de mayor tamaño. La elasticidad de $+0{,}879$ estimada para el Perú es consistente con ese marco: municipios operando por debajo del umbral de escala mínima eficiente exhiben retornos cercanos a la proporcionalidad. \citet{Kinnaman2006} argumenta, desde la perspectiva de la economía del bienestar, que la subprovisión de servicios de recolección en municipios de bajos ingresos no refleja preferencias sino restricciones presupuestales; los resultados de esta tesis respaldan esa lectura para el caso peruano: la elasticidad gasto-QPRS cercana a uno ($+0{,}879$) indica que la capacidad de recolección es aproximadamente proporcional al gasto, por lo que la brecha de cobertura observada entre municipios ricos y pobres es atribuible principalmente a la brecha fiscal y no a diferencias de eficiencia técnica.

La elasticidad de $+0{,}879$ se sitúa en el rango medio-alto de la evidencia comparada en países de ingreso medio-bajo (LMICs).
\citet{Guerrero2013}, en un estudio comparado de sistemas de gestión de residuos en ciudades de América Latina, África Subsahariana y Asia, documentan que la restricción primaria para la recolección formal en municipios de baja capacidad fiscal es el financiamiento operativo, no la disponibilidad tecnológica ni el marco regulatorio.
En los contextos subsaharianos y del Sur de Asia analizados por esos autores, las ciudades con transferencias fiscales más estables exhiben tasas de cobertura de recolección entre 15 y 30 puntos porcentuales superiores a las de ciudades comparables con financiamiento volátil, a igualdad de ingreso per cápita.
Este patrón es consistente con la evidencia sistemática del informe \textit{What a Waste 2.0} \citep{Kaza2018}: en las 368 ciudades de países de ingreso bajo y medio-bajo analizadas, la cobertura de recolección está fuertemente correlacionada con el gasto operativo municipal per cápita, con una pendiente aproximada de $0{,}7$--$0{,}9$ en escala log-log para el cuartil de ciudades más pobres ---rango en el que se ubica la mayor parte de los municipios distritales peruanos.
La elasticidad peruana de $+0{,}879$ es consistente con ese rango y refleja dos características estructurales del contexto.
Primero, la práctica ausencia de provisión informal organizada en los municipios distritales más pequeños implica que el gasto municipal no enfrenta el efecto de sustitución que deprime las elasticidades en contextos con provisión mixta formal-informal.
Segundo, la concentración de municipios con bajos niveles de gasto base ---el $60\%$ gasta menos de S/\,150{,}000 anuales en limpieza--- implica que la mayoría opera en la región de rendimientos crecientes de la función de producción de residuos, donde la elasticidad es necesariamente más alta.
En síntesis, la elasticidad peruana no es idiosincrásica: es la magnitud predecible para municipios de baja capacidad fiscal que operan por debajo del umbral de escala mínima eficiente, donde la restricción es presupuestal y no técnica.

El efecto marginal promedio sobre la probabilidad de uso de relleno sanitario (AME $= +0{,}022$ por incremento del $10\%$ en gasto) es modesto en escala individual pero acumulativamente relevante en el agregado nacional. Proyectado sobre los $1{.}874$ municipios de la muestra, un incremento sostenido del $10\%$ en el gasto de limpieza pública implicaría que aproximadamente $41$ municipios adicionales transitarían hacia disposición formal en un año tipo. La magnitud del AME es comparable con los efectos de subsidios directos a infraestructura de disposición documentados en economías comparables: \citet{Kinnaman2006} reporta que los programas de subsidio a vertederos controlados en jurisdicciones con presupuestos ajustados produjeron transiciones de disposición formal del orden de $1{,}5$ a $3$ puntos porcentuales por ciclo presupuestal. El AME de $+0{,}022$ del Perú, obtenido sin subsidio explícito a la infraestructura sino solo a través del gasto operativo municipal, es comparable con esa magnitud y refuerza la hipótesis de que el financiamiento es la restricción vinculante para la formalización ambiental en municipios de baja capacidad fiscal.

% ------------------------------------------------------------------
\subsection{Validez Externa}
\label{sec:validez_externa}

Los resultados de esta tesis son generalizables a otros contextos de descentralización fiscal en América Latina con estructuras institucionales comparables.
\citet{Faguet2004}, en su análisis del proceso de descentralización boliviano posterior a la Ley de Participación Popular de 1994, documenta que la transferencia de competencias a gobiernos locales aumentó la inversión pública en servicios básicos ---educación, salud, saneamiento--- de forma más sensible a las necesidades locales que el gasto centralizado previo.
El mecanismo identificado por \citet{Faguet2004} es análogo al documentado en esta tesis: las municipalidades con mayor acceso a recursos fiscales reasignan gasto hacia los servicios con mayor déficit relativo, y la restricción vinculante para esta reasignación es financiera, no informacional.
El caso peruano extiende ese hallazgo al servicio de limpieza pública: la elasticidad cercana a la unidad ($+0{,}879$) indica que la capacidad operativa escala casi proporcionalmente con el gasto, lo que implica que la brecha de cobertura entre municipios ricos y pobres es, fundamentalmente, una brecha fiscal transferible mediante la política de transferencias intergubernamentales.

La comparación con países de ingreso bajo y medio fuera de América Latina sugiere que los resultados no son idiosincrásicos al contexto peruano.
\citet{Guerrero2013}, en un análisis comparado de sistemas de gestión de residuos en ciudades de países en desarrollo de América Latina, África Subsahariana y Asia, documentan que la restricción primaria para la formalización de la disposición final es el financiamiento operativo municipal, no la disponibilidad de tecnología ni la regulación.
En los contextos Sub-Saharianos y del Sur de Asia analizados, las ciudades con transferencias fiscales más estables exhiben tasas de disposición formal entre 15 y 30 puntos porcentuales superiores a las de ciudades comparables con financiamiento volátil, a igualdad de nivel de ingreso per cápita.
El AME de $+0{,}022$ estimado para el Perú se sitúa en el rango bajo de este intervalo, lo que es consistente con el hecho de que las transferencias FONCOMUN son relativamente estables en términos reales pero insuficientes en nivel absoluto para la mayoría de municipios distritales rurales.
La convergencia de evidencia entre contextos tan distintos refuerza la interpretación del financiamiento como la restricción vinculante universal para la formalización ambiental en países de ingreso medio y bajo.

% ------------------------------------------------------------------
\subsection{Mecanismos de Transmisión}
\label{sec:mecanismos}

Los resultados son consistentes con dos canales de transmisión distintos, que operan en paralelo pero con horizontes temporales diferentes. El canal operativo directo conecta el gasto con la frecuencia de recolección (FR) y la cantidad recolectada (QPRS): el financiamiento de insumos corrientes ---combustible, mano de obra, mantenimiento de flota--- determina la capacidad operativa inmediata del servicio. Este canal explica por qué FR y QPRS exhiben los efectos más grandes y de más rápida materialización en el panel: son los únicos outcomes que responden al gasto del período corriente sin requerir un proceso burocrático o de inversión intermediario.

El canal institucional opera a través de la planificación. El gasto habilita la contratación de consultores, la elaboración de diagnósticos y los procesos de participación ciudadana que son precondiciones para la aprobación de instrumentos como el PIGARS. La planificación, a su vez, es una precondición administrativa para acceder a financiamiento de infraestructura de disposición final: en el marco regulatorio peruano (Decreto Legislativo N° 1278 y su Reglamento, DS N° 014-2017-MINAM), los municipios sin PIGARS aprobado no son elegibles para financiamiento MINAM de rellenos sanitarios. Esto genera una cadena de causalidad gasto $\to$ PI $\to$ RS que el modelo estima de forma reducida, sin imponer la estructura mediacional, pero que el patrón de coeficientes sugiere: PI responde al gasto con la mayor magnitud relativa entre todos los outcomes ($\hat{\beta} = +0{,}300$), y RS responde con menor magnitud pero de forma estadísticamente inequívoca (AME $= +0{,}022$). El efecto más débil de RS respecto a PI es consistente con la hipótesis de que la planificación es condición necesaria pero no suficiente para la adopción de infraestructura formal: se requieren además inversiones físicas de mayor escala y plazos más largos.

La fortaleza del efecto sobre PI ($\hat{\beta}_{\text{CRE+CF}} = 0{,}300$ vs.\ $\hat{\beta}_{\text{TWFE}} = 0{,}039$, una razón de $7{,}7\times$) refleja que el canal institucional es el más severamente distorsionado por la endogeneidad. Los municipios con mayor capacidad gerencial tienden simultáneamente a ejecutar más presupuesto y a contar con más instrumentos de planificación, de modo que el TWFE sin corrección captura esa covarianza de selección como si fuera efecto del gasto. Una vez corregida la endogeneidad, el efecto puro del financiamiento sobre la planificación es sustancialmente mayor, lo que indica que el gasto es, en sí mismo, un habilitador de la capacidad institucional y no solo un correlato de ella.

% ------------------------------------------------------------------
\subsection{Implicaciones de Política Pública}
\label{sec:politica}

La heterogeneidad implícita en la elasticidad cercana a la unidad ($\hat{\beta}_{\text{QPRS}} = 0{,}879 < 1$) tiene una consecuencia distributiva directa para la asignación presupuestal: los rendimientos marginales del gasto en limpieza pública son decrecientes, por lo que la reasignación inframarginal de recursos desde municipios con alto nivel de gasto base hacia municipios con bajo nivel produciría ganancias netas en la gestión agregada de residuos. Esta es una prescripción de eficiencia dentro del sistema de transferencias existente, sin requerir incremento del gasto total.

El FONCOMUN ---principal transferencia intergubernamental del Perú, distribuida por fórmula basada en población, pobreza y ruralidad--- es el instrumento fiscal más relevante para implementar esa reasignación. Dado que los municipios distritales con menor capacidad fiscal dependen del FONCOMUN para el $80$--$95\%$ de sus ingresos, una modificación de los coeficientes de la fórmula que ponderara el déficit en gestión de residuos sólidos (medido por el índice PI o por la tasa de disposición en botaderos) aumentaría el financiamiento precisamente en los municipios donde el retorno marginal del gasto es mayor. Esta propuesta es consistente con la lógica de asignación por resultados que introduce el DS N° 004-2013-EF sobre presupuesto por resultados en el Perú.

Una estimación de orden de magnitud permite cuantificar el costo fiscal de la reforma.
En el período 2015--2019, el gasto mediano en limpieza pública para los municipios del primer quintil fiscal fue de aproximadamente S/\,85{,}000 anuales, frente a S/\,1{,}200{,}000 para los del quinto quintil (fuente: SIAF/MEF).
Un incremento del $10\%$ en el gasto del primer quintil requeriría S/\,8{,}500 adicionales por municipio-año.
Con un AME de $+0{,}022$, ese incremento produciría una transición de disposición formal en $2{,}2$ de cada $100$ municipios del quintil inferior.
Extrapolado a los $375$ municipios del primer quintil, el incremento generaría aproximadamente $8$ transiciones adicionales por año a un costo fiscal total de S/\,3{,}2 millones anuales ---equivalente a $0{,}006\%$ del presupuesto del Ministerio de Economía y Finanzas en el mismo período.
Para el conjunto de los tres quintiles inferiores (donde el retorno marginal del gasto es mayor según la elasticidad de $+0{,}879$), un incremento del $10\%$ en el gasto de limpieza requeriría un esfuerzo fiscal adicional de S/\,28{,}4 millones anuales, produciendo aproximadamente $41$ transiciones hacia disposición formal ---precisamente el número proyectado a nivel nacional en la \cref{sec:benchmarks}.
La modificación de los coeficientes de la fórmula FONCOMUN para ponderar el déficit de disposición formal requeriría, por tanto, un monto que representa menos del $0{,}05\%$ del presupuesto total de transferencias intergubernamentales, lo que la ubica dentro del rango de reformas fiscalmente neutras implementables sin aumento del gasto total.

Sin embargo, el incremento del gasto sin reforma institucional paralela enfrenta límites estructurales. El AME sobre RS ($+0{,}022$ por $10\%$ de incremento) implica que llevar a un municipio desde botadero a relleno sanitario requiere incrementos del gasto que exceden la capacidad fiscal de la mayoría de municipios distritales rurales bajo el esquema actual. La infraestructura de disposición final requiere inversión de capital con economías de escala que no se logran a nivel individual distrital: los rellenos sanitarios mancomunados, previstos en la Ley N° 29029 (Ley de la Mancomunidad Municipal) y los lineamientos del MINAM, son la solución de política que los datos respaldan. El efecto positivo sobre PRS (AME $= +0{,}026$) en municipios que ya acceden a relleno sanitario indica que el problema de uso insuficiente de infraestructura existente también persiste, y que puede atenderse con incrementos de gasto operativo sin nueva inversión de capital.

% ------------------------------------------------------------------
\subsection{Limitaciones}
\label{sec:limitaciones}

El período de análisis (2015--2019) es previo a la pandemia de COVID-19, que alteró de forma discontinua tanto los patrones de generación de residuos sólidos como los presupuestos municipales en el Perú. Los efectos estimados no son extrapolables al período post-2020 sin verificar que la estructura de la relación gasto-gestión no fue modificada por ese shock. La comparación con datos RENAMU 2020--2023, cuando estén disponibles y auditados por el INEI, permitirá evaluar la estabilidad temporal de los parámetros.

El estimador CRE+CF implementado en esta tesis corrige la endogeneidad bajo el supuesto de que la función de control $\hat{v}_{it}$ captura adecuadamente la correlación entre el gasto y los no-observables. En ausencia de un instrumento externo válido ---la transferencia FONCOMUN requiere el merge con datos MEF/SIAF aún pendiente--- los parámetros estimados no son un estimador de efectos locales promedio (LATE) en sentido estricto. Para cuantificar la sensibilidad a selección no-observable residual, se calculan los límites de plausibilidad de \citet{Oster2019}: el estadístico $\delta^*$ indica cuánto más influyentes que los observables tendrían que ser los no-observables para anular los efectos estimados. Los tres $\delta^*$ son negativos y en valor absoluto superiores a $1{,}3$ (Apéndice~\ref{sec:appendix_oster}), lo que implica que la selección no-observable tendría que operar en dirección opuesta a la selección observable ---condición contraria a la dirección del sesgo documentada por el test de CF--- para reducir los efectos a cero. Este resultado fortalece la plausibilidad de los estimados CRE+CF como efectos condicionales robustos al sesgo de variable omitida estática bajo el supuesto CRE.

Las variables de gestión provienen del Registro Nacional de Municipalidades (RENAMU), que es una declaración jurada sin auditoría física independiente. La calidad del dato depende de la veracidad de los funcionarios municipales en el llenado de los formularios INEI. Subregistro o sobreregistro sistemático de los indicadores de gestión ---especialmente para variables de disposición final, donde la brecha entre lo declarado y lo observado puede ser amplia--- introduciría errores de medición en las variables dependientes, que en presencia de error clásico sesgaría los coeficientes hacia cero (atenuación adicional). En la dirección opuesta, si los municipios con mayor gasto declaran mejores indicadores de forma aspiracional (informe de la gestión presupuestal), el error de medición generaría sesgo al alza. Ambas posibilidades son plausibles y no descartables sin un estudio de validación con datos de campo.

Los controles socioeconómicos externos ---pobreza provincial (ENAHO), urbanización y población (CPV 2017)--- fueron interpolados linealmente entre los años de levantamiento de las encuestas fuente, introduciendo error de medición en los regresores de control. Dado que estos controles no son la variable de interés, el efecto sobre los estimados de $\hat{\beta}$ depende del grado de correlación entre el gasto y los controles interpolados; el análisis de sensibilidad reportado en el Apéndice sugiere que la dirección y magnitud de los coeficientes principales es estable ante variaciones en los controles.

% ------------------------------------------------------------------
\subsection{Agenda de Investigación}
\label{sec:agenda}

La extensión más inmediata y de mayor valor para la identificación es el merge del gasto SIAF/MEF con los datos FONCOMUN del mismo período. Con esa base, la fórmula de distribución del FONCOMUN ---cuyos coeficientes son fijados centralmente y no responden a decisiones de los municipios individuales--- puede usarse como instrumento externo para $\ln G_{it}$, transformando el estimador CRE+CF en un IV-2SLS dentro del marco CRE. El estimador IV recuperaría el LATE para el subgrupo de municipios cuyo gasto en limpieza pública responde a variaciones en la transferencia FONCOMUN ---previsiblemente los municipios de menor capacidad fiscal y mayor dependencia de transferencias, que son precisamente aquellos con mayor retorno marginal del gasto según la elasticidad estimada.

El análisis de heterogeneidad por tamaño de municipio y región geográfica constituye una segunda extensión prioritaria. La elasticidad agregada de $0{,}879$ oculta una distribución de efectos que puede variar sustancialmente entre municipios urbanos grandes (donde los problemas de gestión son de escala y logística) y municipios rurales pequeños (donde la restricción es de capacidad operativa básica). Una especificación con interacciones entre $\ln G_{it}$ y cuartiles de población, o entre $\ln G_{it}$ y dummies de región natural (costa, sierra, selva), identificaría si los rendimientos decrecientes del gasto operan de forma homogénea o si existen subpoblaciones con elasticidades significativamente más altas donde la intervención fiscal tendría mayor impacto.

La extensión temporal del panel hacia los años 2020--2023, condicionada a la disponibilidad del RENAMU correspondiente, permitiría dos análisis adicionales. Primero, evaluar si el shock COVID-19 generó un quiebre estructural en la relación gasto-gestión, lo que tiene implicaciones para la validez externa de los parámetros estimados. Segundo, explotar la variación en los programas de incentivos fiscales municipales (Plan de Incentivos a la Mejora de la Gestión Municipal) que el MEF ha modificado en ese período, generando una fuente adicional de variación cuasi-experimental en el gasto de limpieza.
